\documentclass[border=8pt, multi, tikz]{standalone} 
\usepackage{import}
\subimport{../layers/}{init}
\usetikzlibrary{positioning}
\usetikzlibrary{calc}
\usetikzlibrary{3d} %for including external image 

\def\ConvColor{rgb:yellow,5;red,2.5;white,5}
\def\BNColor{rgb:blue,2;green,1;white,4}
\def\ReLUColor{rgb:orange,5;white,3;black,2}
\def\GAPColor{rgb:blue,2;black,1;white,4}
\def\ResBlockColor{rgb:blue,5;red,2.5;white,5}
\def\ConvReluColor{rgb:yellow,5;red,5;white,5}
\def\PoolColor{rgb:red,1;black,0.3}
\def\UnpoolColor{rgb:blue,2;green,1;black,0.3}
\def\FcColor{rgb:blue,5;red,2.5;white,5}
\def\FcReluColor{rgb:blue,5;red,5;white,4}
\def\SoftmaxColor{rgb:magenta,5;black,7}   
\def\SumColor{rgb:blue,5;green,15}

% colour of the shortcut branch, reused for the branch dot
\def\SkipColor{rgb:blue,4;red,1;green,1;black,3}

\newcommand{\copymidarrow}{\tikz \draw[-Stealth,line width=0.8mm,draw=\SkipColor] (-0.3,0) -- ++(0.3,0);}

\begin{document}
\begin{tikzpicture}
\tikzstyle{connection}=[ultra thick,every node/.style={sloped,allow upside down},draw=\edgecolor,opacity=0.7]
\tikzstyle{copyconnection}=[ultra thick,every node/.style={sloped,allow upside down},draw=\SkipColor,opacity=0.7]

% height of the horizontal run of the shortcut
\def\skipy{4.9}

%%%%%%%%%%%%%%%%%%%%%%%%%%%%%%%%%%%%%%%%%%%%%%%%%%%%%%%%%%%%%%%%
%  main branch:  Conv2D -> BN -> ReLU -> Conv2D -> BN -> (+)
%%%%%%%%%%%%%%%%%%%%%%%%%%%%%%%%%%%%%%%%%%%%%%%%%%%%%%%%%%%%%%%%
\pic[shift={ (0,0,0) }] at (0,0,0) { 
    Box={
        name=conv1,
        caption=Conv2D,
        fill=\ConvColor,
        height=32,
        width=2,
        depth=32
        }
    };

\pic[shift={ (1.2,0,0) }] at (conv1-east) { 
    Box={
        name=bn1,
        caption=BN,
        xlabel={ {  }, },
        zlabel={},
        fill=\BNColor,
        height=32,
        width=0.4,
        depth=32
    }
};
\draw [connection]  (conv1-east)    -- node {\midarrow} (bn1-west);

\pic[shift={ (1.5,0,0) }] at (bn1-east) {
    Box={
        name=relu1,
        caption=ReLU,
        xlabel={ {  }, },
        zlabel={},
        fill=\FcReluColor,
        height=32,
        width=0.4,
        depth=32
    }
};
\draw [connection]  (bn1-east)    -- node {\midarrow} (relu1-west);

\pic[shift={ (1.5,0,0) }] at (relu1-east) { 
    Box={
        name=conv2,
        caption=Conv2D,
        fill=\ConvColor,
        height=32,
        width=2,
        depth=32
        }
    };
\draw [connection]  (relu1-east)    -- node {\midarrow} (conv2-west);

\pic[shift={ (1.2,0,0) }] at (conv2-east) { 
    Box={
        name=bn2,
        caption=BN,
        xlabel={ {  }, },
        zlabel={},
        fill=\BNColor,
        height=32,
        width=0.4,
        depth=32
    }
};
\draw [connection]  (conv2-east)    -- node {\midarrow} (bn2-west);

\pic[shift={(2.4,0,0)}] at (bn2-east) 
    {Ball={
        name=sum1,
        fill=\SumColor,
        opacity=0.6,
        radius=2.5,
        logo=$+$
        }
    };
\draw [connection]  (bn2-east)    -- node {\midarrow} (sum1-west);

%%%%%%%%%%%%%%%%%%%%%%%%%%%%%%%%%%%%%%%%%%%%%%%%%%%%%%%%%%%%%%%%
%  block input, block output
%%%%%%%%%%%%%%%%%%%%%%%%%%%%%%%%%%%%%%%%%%%%%%%%%%%%%%%%%%%%%%%%
\coordinate (blockin)  at (-4.0,0,0);
\coordinate (branch)   at (-2.4,0,0);   % where the shortcut leaves the input
\coordinate (blockout) at ($(sum1-east)+(2.6,0,0)$);

% input -> first Conv2D
\draw [connection] (blockin) -- node {\midarrow} (conv1-west);

% sum -> block output
\draw [connection] (sum1-east) -- node {\midarrow} (blockout);

\node[anchor=east, font=\small] at ($(blockin)+(-0.15,0,0)$)  {Input};
\node[anchor=west, font=\small] at ($(blockout)+(0.15,0,0)$)  {Output};

%%%%%%%%%%%%%%%%%%%%%%%%%%%%%%%%%%%%%%%%%%%%%%%%%%%%%%%%%%%%%%%%
%  identity shortcut: taken from the INPUT of the block
%%%%%%%%%%%%%%%%%%%%%%%%%%%%%%%%%%%%%%%%%%%%%%%%%%%%%%%%%%%%%%%%
\coordinate (skipA) at ($(branch)+(0,\skipy,0)$);
\coordinate (skipB) at (skipA -| sum1-north);

\draw [copyconnection]
      (branch)
  --  node {\copymidarrow} (skipA)
  --  node {\copymidarrow} (skipB)
  --  node {\copymidarrow} (sum1-north);

% dot marking the branch point on the input line
\fill[draw=\SkipColor, fill=\SkipColor] (branch) circle (3.2pt);

\node[anchor=south, font=\small] at ($(skipA)!0.5!(skipB)+(0,0.30,0)$)
  {identity shortcut};

\end{tikzpicture}
\end{document}